\pagebreak

{
\thispagestyle{empty}

\mbox{}

\vspace{11cm}

\footnotesize
\noindent{}\emph{\noindent{}``Um \emph{hassid} perguntou ao Maguid de Zlotschov. Nosso mestre Rashi diz: O que faltava ao mundo adormecido? Só o descanso. Veio o \emph{shabat} e, com ele, o descanso. Por que não se diz simplesmente que ao mundo faltava o descanso, até que veio o \emph{shabat}? Pois se \emph{shabat} e descanso querem dizer a mesma coisa!\medskip\\
--- \emph{Shabat}, disse o Rabi, quer dizer a volta ao lar, porque neste
dia todas as esferas voltam ao seu verdadeiro lugar. É para isto que
aponta Rashi: inquietas estão as esferas durante a semana, porque foram
descidas de seus lugares, e no \emph{shabat} encontram o descanso, porque podem
voltar.''\footnote{\textsc{buber}, \emph{op.\,cit.}, p.\,196.}}
}

\pagebreak

\chapter{Geografias}

% \begin{quote}
% ``Um \emph{hassid} perguntou ao Maguid de Zlotschov. Nosso mestre Rashi diz: \emph{O
% que faltava ao mundo adormecido? Só o descanso. Veio o \emph{shabat} e, com
% ele, o descanso}. Por que não se diz simplesmente que ao mundo faltava o
% descanso, até que veio o \emph{shabat}? Pois se \emph{shabat} e descanso querem dizer
% a mesma coisa!

% --- \emph{Shabat} --- disse o Rabi --- quer dizer a volta ao lar, porque neste
% dia todas as esferas voltam ao seu verdadeiro lugar. É para isto que
% aponta Rashi: \emph{inquietas estão as esferas durante a semana, porque foram
% descidas de seus lugares, e no \emph{shabat} encontram o descanso, porque podem
% voltar}''\footnote{\textsc{buber}, \emph{op.\,cit.}, p.\,196.}.
% \end{quote}

\noindent{}Segundo a filosofia judaica, o \emph{shabat} é o sétimo dia da criação
divina, momento em que a divindade apenas contemplou o mundo e suas
coisas. Como recordação da criação do mundo por Deus, o \emph{shabat} é
considerado um período sagrado e dedicado exclusivamente ao repouso.
Inicia-se no pôr do sol da sexta-feira e termina ao nascer das três
primeiras estrelas do sábado, totalizando 25 horas de duração.

Durante o \emph{shabat}, judeus e judias recordam o ato de criação
divina do mundo --- e de todas as coisas que ele contém --- por meio do
descanso físico, mental e espiritual e da pausa do trabalho criativo. O
momento do \emph{shabat}, dessa forma, se mostra enquanto um espaço (e
um tempo) em que se permite o ócio, propício a reflexões, ``uma
volta ao lar'', segundo o Maguid de Zlotschov --- e assim a conexão com
a divindade e com a comunidade\footnote{As comunidades judaicas se unem
  e se apropriam de um território real para criarem seu mundo e
  realizarem seus ritos religiosos.}. Afinal, ``sem o ócio não
poderia haver criação, apenas repetição monótona do mesmo, submissão à
esfera da necessidade''\footnote{\textsc{rago}, \emph{op.\,cit.}, p.\,10.}.

Nessa lógica, são ao todo 39 \emph{melachot}\footnote{\emph{Melachot}
  (plural de \emph{melachá}) são ações criativas que interferem
  diretamente na lei do descanso judaico durante o \emph{shabat}. As
  \emph{melachot} são atividades que alteram a ordem natural das coisas,
  bem como a maneira como se encontrava o mundo no exato momento em que
  foi concebido o \emph{shabat}.} proibidas durante o rito do
\emph{shabat}, sendo elas: plantar, arar, cortar, juntar, debulhar,
dispersar, separar, moer, peneirar, fazer massa, assar, tosquiar, lavar,
desembaraçar, tingir, fiar, esticar, passar, tecer, desfazer nós e/\,ou
fios, desatar, costurar, rasgar, caçar, abater, pelar o couro, curtir o
couro, demarcar o couro, cortar, escrever, apagar, construir, destruir,
apagar o fogo, acender o fogo, terminar a manufatura de um objeto,
transportar de uma propriedade à outra ou carregar em um domínio
público.

Durante o \emph{shabat}, carregar objetos ou corpos, entre domínios
públicos\footnote{Segundo as interpretações rabínicas, domínios públicos
  ou, em hebraico, \emph{reshut harabin}, são todos aqueles espaços-ruas
  que conectam uma cidade de uma ponta a outra, que tenham pelo menos
  doze metros de largura e onde transitem 600 mil pessoas por dia. Um
  exemplo de domínio público na cidade de São Paulo seria a avenida
  Paulista.} e domínios privados\footnote{Segundo as interpretações
  rabínicas, domínios privados ou, em hebraico, \emph{reshut hayachid},
  são todos aqueles espaços que tenham pelo menos 80 cm de elevação do
  solo e uma área mínima de 4 por 4 \emph{tefachim}, aproximadamente 40
  x 40 cm, podendo um indivíduo e/\,ou associação ser proprietária desse
  local, ou não.} (e vice-versa), é proibido segundo a
\emph{Halachá}\footnote{\emph{Halachá} é um conjunto de normas e leis
  que servem como guia para uma vida judaica religiosa. A palavra
  \emph{Halachá} vem do vocabulário hebraico, mais especificamente do
  termo \emph{heh-lamed-kaf}, que significa viagem e/\,ou caminho.}.
Assim, é dentro dessa perspectiva que surge o \emph{eruv}. Um artifício
que se estabelece enquanto um portal nas territorialidades em que as
comunidades se agregam, ou seja, ``uma linha que permite a junção
de vários tempos e espaços, e que ali, neste ponto para o qual tudo
converge, seria possível o trânsito entre passado e futuro, sonho e
vigília, vida e morte''\footnote{\textsc{saavedra}, \emph{O manto da noite}, p.\,86.}.

O \emph{eruv}, por assim dizer, é um enclausuramento ritualístico que
permite a abertura dos espaços privados da morada ao domínio público ---
criando, desta forma, uma casa aberta com becos, vielas e pátios
compartilhados entre a comunidade --- permitindo a circulação e o
transporte de objetos.

A fantasia do \emph{eruv} acontece às escondidas. Este território mágico
surge nas frestas e nas brechas de uma proibição judaica religiosa --- o
carregamento de objetos durante o rito judaico do \emph{shabat}. O
\emph{eruv} materializa o texto --- e suas interpretações --- no espaço
e se justifica a partir dos debates rabínicos sobre os livros sagrados.
O \emph{eruv} é uma estratégia que contesta as restrições\footnote{Mas a
  partir das regras do jogo --- ou seja, a partir das interpretações
  rabínicas dos textos.} de carregamento, e mais que nada entende que,
para além das práticas religiosas, as comunidades judaicas sobrevivem em
diásporas a partir de suas agregações, seus \emph{minianim}\footnote{\emph{Minianim}
  (plural de \emph{minian}) refere-se ao quórum mínimo de dez homens
  judeus adultos necessários para realizar certas obrigações religiosas,
  como a retirada da Torá da Arca Sagrada e sua leitura. Mais
  recentemente e em algumas vertentes religiosas e liberais judaicas, as
  mulheres também passaram a contabilizar como participantes ativas do
  \emph{minian}.}, seus laços e traços de pertencimento e de
sociabilidade.

Para a pesquisadora e arquiteta Miriam Levy, ``o \emph{eruv} é uma forma
de navegação nômade sobre a topografia de um espaço, em que a noção de
pertencimento a uma comunidade é criada e construída em detrimento de um
território em particular''\footnote{\textsc{levy}, \emph{Encounter with the
  \emph{eruv}: a project toward the city of open enclosures}, p.\,51. (tradução
  nossa)}. Por fim, independentemente deste local, o \emph{eruv} pode
ser acionado e construído em qualquer lugar e em qualquer tempo, desde
que haja uma comunidade agregada.

\begin{center}\adforn{49}\end{center}

\noindent{}\emph{Eruv} provém da palavra aramaica \emph{arab}, que em português
significa ``mistura''. Ou seja, o \emph{eruv} é um artifício judaico que
permite a combinação de espaços públicos e espaços privados,
apropriando-se destas territorialidades para criar assim um outro local
de pertencimento, de trânsito e de circulação. De forma prática, este
artifício se estabelece no tecido urbano das cidades por meio de um
\emph{tsurat hapetach} (em português: ``formas de portas''), dois postes
verticais metálicos implementados paralelamente às calçadas, rentes às
construções da cidade, e conectados em suas extremidades por um fio de
nylon transparente ou por um barbante.

A marcação deste espaço físico surge enquanto resposta às restrições de
carregamento no \emph{shabat}. Assim, o \emph{eruv} é nômade e
construído por meio de conexões e interferências com os espaços urbanos
e suas edificações pré-existentes. Em um contexto diaspórico há uma
necessidade de criar e consolidar uma casa comum e compartilhada em um
determinado território --- apropriar-se dele, consagrá-lo desde dentro.
Por meio desta estrutura física e simbólica permite-se a solidificação
de espaços judaicos em que a comunidade se conecta com uma cultura
pré-existente e, simultaneamente, também com a sua própria. E assim,
``o \emph{eruv} faz parte da topografia urbana da cidade e interage com
ela''\footnote{\emph{Ibidem}, p.\,13. Tradução da autora.}.

O \emph{eruv} possibilita o estreitamento de laços dentro da comunidade
que está inserida em um determinado território por ele delimitado. A
apropriação destes espaços molda as dinâmicas sociais, religiosas,
simbólicas e seus deslocamentos durante o rito do \emph{shabat}.
Entende-se aqui que o \emph{eruv} é um artifício de pertencimento por
meio do qual indivíduos judeus se locomovem nos territórios em que se
fixam.

Estes deslocamentos entre espaços públicos e espaços privados permitem
uma identificação e uma agregação entre os membros da comunidade
judaica; ao final, o contato visual nas ruas entre minorias e com uma
cultura maior (sociedade) estreita o senso de comunidade, de
identificação, e ajuda os indivíduos a formar laços políticos e alianças
sociais de resistência nos territórios diaspóricos.

Dessa forma, ao sair às ruas, a comunidade torna o ambiente externo
menos ameaçador, faz dele uma localidade apropriada. Ao final, o senso
de pertencimento e as agregações comunitárias são atos de resistência
que demarcam e delimitam territórios, solidificam as relações
interpessoais e apresentam a comunidade enquanto um coletivo.

\begin{center}\adforn{49}\end{center}

\noindent{}O \emph{eruv} é uma área circunscrita por meio da união de construções
pré-existentes e de fios --- paredes físicas e paredes simbólicas,
respectivamente --- em que se permite o enclausuramento e a apropriação
destas territorialidades. É por meio da apropriação deste local que se
abre uma ruptura no espaço-tempo, possibilitando relações de
sociabilidade durante o rito do \emph{shabat}. Assim, como forma de
cumprir as restrições de carregamento do \emph{shabat}, os espaços
urbanos, supostamente laicizados, são alterados e ressignificados pelas
ações e pelas interações e interpretações religiosas.

Neste sentido, o \emph{eruv} se torna uma casa-território compartilhada
pela comunidade judaica local durante o rito do \emph{shabat}. É uma
territorialidade em que os espaços coexistem, porém não mais coordenados
pelo binômio público-privado\footnote{O Talmud divide os espaços
  e as territorialidades em duas categorias: \emph{reshut harabin} e
  \emph{reshut hayachid} (espaços públicos e espaços privados,
  respectivamente). Os rabinos, a partir das interpretações das
  características espaciais, definem mais dois lugares, \emph{carmelit}
  e espaços neutros. Os espaços \emph{carmelit} são locais que
  apresentam características de espaços públicos, porém sem as dimensões
  exatas necessárias, por exemplo, rua de bairro, vila de casas. Já os
  espaços neutros são territórios que se assemelham aos espaços
  privados, mas que possuem menos de 1,6m² de área ou 80 cm de altura,
  por exemplo, postes de eletricidade, decks de piscinas, tábua de uma
  mesa.}. Segundo Miriam Levy, ``é nestes locais que é possível
encontrar outros sinais de identidade e novos espaços de
colaboração''\footnote{\textsc{levy}, \emph{op.\,cit.}, p.\,21. Tradução da autora.}.

O \emph{eruv} é uma mescla de pátios. Este encontro de domínios, espaços
públicos e espaços privados apropriados se dá por meio de uma refeição
compartilhada entre a comunidade. Pois, no caso judaico, a comunidade é
um marcador legal e social de pertencimento. Todas as pessoas que ali se
alimentam fazem parte de uma mesma comunidade\footnote{Os primeiros
  \emph{eruvim} (plural de \emph{eruv}) surgem nas comunidades judaicas
  logo após a destruição do segundo Templo Sagrado de Jerusalém (76
  d.E.C.), como uma tentativa de criação de comunidade para além das
  restrições do \emph{shabat}. Este convívio comunitário era realizado
  por meio de uma associação entre os diversos moradores das casas, que
  compunham uma vila, por meio de uma refeição compartilhada. No momento
  do \emph{shabat}, as vilas judaicas eram cercadas por tábuas de
  madeira posicionadas nos topos das construções, e assim simbolicamente
  criavam um umbral/\,porta de pertencimento à localidade.}. Essa
demarcação constrói uma casa simbólica no território da cidade, e é
entoada por meio da bênção do \emph{eruv}:

\begin{quote}
Este é o \emph{eruv} que deve ser utilizado para permitir que nós possamos
trazer coisas de dentro para fora. De nossas casas para os nossos
quintais. De nossos quintais para nossas casas. De casa em casa, de
quintal em quintal, de terraço em terraço. De nossas casas e nossos
quintais para becos e vielas, e de becos e vielas para suas casas. Para
nós, e para todas as pessoas que vivem aqui, e quem se juntar conosco
durante o rito do \emph{shabat} e nas festividades\footnote{Bênção judaica do
  \emph{eruv} entoada sobre um alimento que será repartido entre a
  comunidade.}.
\end{quote}

``A comida'', como esclarece Levy, ``é um símbolo de que
não importa em que localidade estas pessoas estejam; porém, entende-se
que esses indivíduos estão em um determinado espaço e tempo em que está
também a comunidade, e assim, onde a comunidade se
estabelece''\footnote{\textsc{levy}, \emph{op.\,cit.}, p.\,59. Tradução da autora.}. Na
maior parte dos casos, essa comida compartilhada em comunidade é uma
\emph{matzá}, um pão ázimo sem fermentação, como preparado originalmente
durante a fuga bíblica da escravidão no Egito. Simboliza, ao mesmo
tempo, a transitoriedade --- o pão que não ficou pronto --- e o
pertencimento.

É a \emph{matzá} --- essa refeição --- que é expandida simbolicamente no
território físico e urbano da cidade, sendo ela também o próprio
\emph{eruv} --- possibilitando que toda a comunidade compartilhe este
alimento. Os fios e postes, ao final, são delimitadores físicos (pouco
visíveis, se não invisíveis para muitos) no território urbano --- a
partir deles entende-se quais ruas estão englobadas e quais ruas estão
descartadas dentro da marcação do ``Distrito de
\emph{Eruv}''\footnote{Apesar do \emph{eruv}, em grande maioria, se
  apresentar enquanto uma construção fictícia, imaginária e silenciosa
  de bairro e de agregação comunitária, a consolidação do \emph{eruv} na
  cidade de São Paulo necessitou de um contrato legal e jurídico entre a
  comunidade judaica e a prefeitura --- esta permite que a comunidade
  alugue simbolicamente os direitos do Tesouro Municipal, para, dessa
  forma, implementar os postes nas calçadas públicas da cidade. O
  contrato, assim sendo, garante que a comunidade judaica seja
  legalmente (e simbolicamente) proprietária do território no qual o
  \emph{eruv} é instalado --- chamado de ``Distrito de \emph{Eruv}''.}.

O \emph{eruv} é um marcador legal e social de pertencimento comunitário.
Estar dentro desta territorialidade, uma construção fictícia de bairro
--- inventado e apropriado pelas comunidades judaicas ---, permite uma
simbolização. Um estabelecimento \emph{kosher}\footnote{\emph{Kosher}
  significa ``apropriado'' e/\,ou ``adequado'' e são preceitos e práticas
  que seguem as normas da religiosidade judaica. Um exemplo são os
  alimentos \emph{kosher}, preparados segundo as leis alimentares
  judaicas, sendo elas: carne de animais que sejam, ao mesmo tempo,
  ruminantes e de patas fendidas, peixes que possuem escamas e
  barbatanas, além da proibição da mistura de carne com laticínios.},
apesar de não funcionar durante o \emph{shabat}, necessita estar dentro
das delimitações do \emph{eruv}. Assim como estão as sinagogas, os
centros comunitários, as escolas e as casas judaicas.

O \emph{eruv}, neste sentido, é um assunto comunitário. Durante o
\emph{shabat}, há a obrigação religiosa da ida à sinagoga. Este espaço,
mais do que uma casa de orações, é um centro comunitário, de convívio,
de estudos --- e assim, de socialização. ``O \emph{eruv} é sobre a
movimentação de pessoas que estão em suas casas para as
ruas''\footnote{\textsc{levy}, \emph{op.\,cit.}, p.\,33. Tradução da autora.}. Se, por conta
das leis restritivas do \emph{shabat}, pessoas com necessidades
especiais não pudessem sair de casa, a comunidade estaria inviabilizando
uma vida judaica coletiva. Sem o \emph{eruv}, indivíduos que utilizam
bengalas, cadeiras de rodas, pessoas com mobilidade reduzida, portadores
de deficiências físicas, adultos com crianças de colo que necessitam do
auxílio de carrinhos de bebês não se deslocariam de suas casas para a
confraternização com a comunidade e, assim, ficariam privados de viver
um judaísmo em coletivo.

Nesta perspectiva, o artifício do \emph{eruv} se sustenta enquanto um
fortificador diaspórico de comunidade, no sentido mais amplo desta
palavra. Afinal, viver em comunidade, segundo o escritor e pesquisador
Luiz Antônio Simas, ``representa a agregação de diversos
indivíduos em busca de objetivos comuns''\footnote{\textsc{simas}; \textsc{rufino}, \emph{op.\,cit.}, p.\,4.}, ou seja, criar e manter conexões --- mais que nada,
entende-se aqui que pessoas com dores, e com sonhos, criam comunidades.

\begin{center}\adforn{49}\end{center}

\noindent{}``O \emph{eruv} é a construção de um único espaço privado/\,domesticado ao
redor de um bairro/\,comunidade, em que carregar objetos torna-se
permitido''\footnote{\textsc{levy}, \emph{op.\,cit.}, p.\,83. Tradução da autora.}. Assim, a
partir de leis religiosas, a estratégia do \emph{eruv} cria uma série de
ressignificações nos espaços físicos e urbanos das cidades
e ``apaga as divisões legais entre os espaços, cria um grande
espaço privado, domesticado''\footnote{\emph{Ibidem}, p.\,36. Tradução da autora.}.

As leis judaicas religiosas do \emph{shabat} permitem o transporte de
objetos e utensílios nos espaços privados, mas não entre um espaço
privado e outro público. Tal interdição justifica-se porque os objetos
--- bolsas, celulares, carteiras, livros etc. --- podem ser obstáculos a
relações humanas mais aprofundadas. Como nos lembra o Maguid de
Zlotschov, ``o \emph{shabat} é uma volta ao lar''\footnote{\textsc{buber}, op.
  cit., p.\,196.} em que as relações interpessoais são intensificadas.

O \emph{shabat} é um convite para olhar, para refletir e para se
encontrar com o outro dentro de si mesmo. Tudo no \emph{shabat} está
voltado e vinculado ao descanso mental e espiritual. Nesse período,
todos os atos, ritos e práticas judaicas se relacionam ao conceito de
\emph{shalom} (paz).\footnote{Em ``O Santo \emph{Shabat}'': ``A
  mãe do Rabi Mendel perguntou-lhe certa vez: \emph{O que significa realmente
  dizermos: o santo \emph{shabat}?} Ele respondeu: \emph{O \emph{shabat} santifica}. Ela,
  porém, disse: \emph{Não só santifica, como também cura}'. \textsc{buber}, \emph{op.\,cit.}, p.\,446.} Este encontro com a paz parte do
convívio comunitário, no entrelaçamento entre a comunidade, o descanso,
a religiosidade e a divindade. Afinal, os seres humanos foram criados
para terem prazer com a conexão divina.

A crença central do \emph{shabat} procura desenvolver um espírito de
paz. O \emph{shabat} começa com a imagem de uma noiva e de um noivo
caminhando em direção à \emph{chupá}\footnote{\emph{Chupá} é uma tenda
  sob a qual se realiza o casamento judaico.} e termina com a presença
do profeta Eliahu, anunciador da era messiânica e da vinda do terceiro
Templo Sagrado de Jerusalém. A união dos noivos é expandida de forma a
incluir a esperança da redenção do mundo, por completo\footnote{Ao
  final, ``Aquele que estabelece a paz nas alturas, há de
  estabelecer a paz entre nós''. Jó 25:2}. O \emph{shabat} começa e termina com acendimento de
fogo. Em um primeiro momento, ao pôr do sol da sexta-feira, utiliza-se
duas velas independentes --- dois castiçais acendidos separadamente. E,
ao final do período, ao nascer das três primeiras estrelas ao sábado,
utiliza-se uma única vela de \emph{havdalá},\footnote{\emph{Havdalá} é
  uma cerimônia religiosa judaica que marca o fim do período do
  \emph{shabat}, dando início a uma nova semana.} em formato de trança
que se enrosca.

\begin{quote}
O \emph{shabat} começa com a oração que descreve o primeiro \emph{shabat} no Éden,
onde existiam apenas Adão e Eva, e termina com a esperança e a oração de
que a unicidade de Deus se reflita na unicidade do povo.\footnote{WEISS,
  \emph{A Microcosm of the Shabbat Spirit}, p.\,5. Tradução da autora.}
\end{quote}

Se a unicidade de Deus está refletida na unicidade do povo, então a
filosofia judaica parte do pressuposto de que já há uma familiaridade e
unicidade (um sentido de paz) entre os diversos indivíduos que
compartilham um determinado espaço de convivência; e assim, faz-se
permitido o transporte de objetos dentro dos espaços privados e/\,ou
domesticados --- a casa. Ao final, a casa, a morada, por si só
estabelece uma pequena comunidade, desde que haja uma certa intimidade
entre os indivíduos que nela habitam.

Assim, ao cercar uma determinada territorialidade da cidade, o
\emph{eruv} não incide diretamente em espaços públicos como os
conhecemos. O \emph{eruv} cerca e delimita espaços urbanos considerados,
segundo as interpretações rabínicas, espaços \emph{carmelit}.
\emph{Carmelit} são as ruas de bairro, pequenas vilas de casas, ruas com
comércios menos movimentados. O \emph{eruv}, ao circunscrever-se nos
espaços \emph{carmelit} --- uma territorialidade de pouco trânsito ---,
pressupõe que os indivíduos que transitam por estas zonas já fazem parte
de um microcosmo --- uma ideia de casa, de intimidade ---, a comunidade
em si já existe\footnote{Este é o caso da Festa do \emph{Eruv},
  realizada pela Casa do Povo anualmente e vista pela comunidade do Bom
  Retiro enquanto uma festa de bairro, apropriada pelos indivíduos que
  transitam, moram ou moram ali. Por assim dizer, mesmo que o conceito
  do \emph{eruv} tenha nascido dentro de um ambiente judaico, a Casa do
  Povo, durante a Festa do \emph{Eruv}, traduz este conceito e engloba
  todas as pessoas, transformando o espaço da Casa em um território
  compartilhado. Assim, ao se apropriar do espaço público da rua Três
  Rios (local onde está situada a Casa do Povo no Bom Retiro), a Festa
  do \emph{Eruv} convida toda a comunidade do bairro a frequentar este
  território. Durante o evento diversos são os indivíduos que participam
  ativamente na organização, nas oficinas e nas atividades propostas.
  Afinal, o Bom Retiro é conhecido como um bairro multicultural, em que
  transitam e onde se estabelecem diversos grupos sociais: judeus,
  coreanos, paraguaios, bolivianos, venezuelanos, nordestinos e
  comunidades negras, pessoas LGBT+, prostitutas, gregos. A casa se
  torna um espaço comum, uma outra frequência.}.

Se partirmos do pressuposto de que o \emph{eruv} combina a arquitetura
local com uma arquitetura judaica religiosa --- um espaço (e um tempo)
físico e imaginário, real e ficcional, concreto e abstrato ---, então
ele constrói para si um ``espaço de pertencimento que é nômade em
sua forma de construção, mas estacionário em suas conotações
metafóricas''\footnote{\textsc{levy}, \emph{op.\,cit.}, p.\,61. Tradução da autora.}. O
\emph{eruv}, paradoxalmente, cria um espaço em que há um sentimento
comunitário de pertencimento sedentário, até que ele se desmanche ---
``este local de pertencimento se move com a comunidade, quando
necessário''\footnote{\emph{Ibidem}, p.\,59. Tradução da autora.} --- e seja
construído em outro território.

\begin{center}\adforn{49}\end{center}

\noindent{}O \emph{eruv} é uma fronteira, um limite e um limiar em torno de toda
uma comunidade e se estabelece enquanto uma multiplicidade de espaços e
de tempos sagrados e multifacetados. Este espaço sagrado (e apropriado)
engloba pessoas, objetos, animais, textos, livros e atividades
ritualísticas, religiosas, cotidianas, coletivas e comuns, bem como as
suas arquiteturas. Durante o \emph{shabat}, a comunidade judaica é
unificada em um espaço de pertencimento e de identificação, tornando
esta territorialidade um espaço sagrado, ``Seu Mundo'', um local
que permite o descanso (e o encontro com a paz) e uma territorialidade
de presença e de devoção divina. Como menciona a escritora bell hooks:

\begin{quote}
Os espaços podem ser reais e imaginários. Os espaços podem contar
histórias e desdobrar relatos históricos. Os espaços podem ser
interrompidos, apropriados e transformados pela prática artística e
literária\footnote{\textsc{hooks}; \textsc{yearning}. \emph{Race, gender and cultural
  politics}, p.\,23. Tradução da autora.}.
\end{quote}

Enquanto um artifício, um objeto e um território temporário, e assim,
diaspórico, o \emph{eruv} não acontece, nem mesmo pertence, a um
determinado espaço, já que ``um espaço se torna um lugar quando se
coloca uma série de objetos nele''\footnote{\textsc{levy}, \emph{op.\,cit.}, p.\,59. Tradução da autora.}. O artifício do \emph{eruv} parte de seu
agenciamento, e a ``agência dos objetos'' permite que estes artifícios
se tornem suporte para algo invisível.

O \emph{eruv} não é um não lugar, ou melhor, o artifício do \emph{eruv}
não se estabelece enquanto uma utopia. O \emph{eruv} é um território
real e físico --- um espaço-tempo apropriado pelas comunidades judaicas
ortodoxas ---, em que outros tempos, outras frequências e outras
relações de sociabilidade se fixam, se concretizam, manifestam-se para
que, então, desmanchem-se\footnote{Desmanchar é o ato de alterar ou
  desfazer o aspecto, a forma ou a arrumação de um determinado objeto,
  bem como desarranjar, desalinhar e desarrumar suas estruturas.
  ``Desmanchar'' é o verbo que permite que partículas moleculares se
  desintegrem. O objeto/\,corpo simplesmente some e seus resquícios ---
  que passam a ser microscópicos --- voam pelo ar. O \emph{eruv}, em seu
  aspecto visível-invisível, se desintegra com facilidade. Por ser um
  objeto móvel e diaspórico permite que seus limites se expandam (ou se
  retraiam), construindo e consolidando localidades --- até que estes
  locais se desvaneçam.}. O \emph{eruv} é uma heterotopia
territorial\footnote{Heterotopia é um conceito da geografia humana que
  descreve lugares que funcionam em condições não hegemônicas, com
  múltiplas camadas de significação. As heterotopias são espaços de
  alteridade, simultaneamente físicos e mentais.} que permite pensar os
espaços (e os tempos) como multiformes --- um interminável e simultâneo
cruzamento de localidades, de tempos e de situações.

Assim, o \emph{eruv} pode ser analisado enquanto um signo da diáspora,
já que se estabelece enquanto uma construção real e fictícia situada em
um determinado espaço urbano em que a comunidade decide fixar-se. O
\emph{eruv} são fios e postes móveis, que se expandem ou se encolhem a
depender das necessidades, dos laços e dos traços de sociabilidade. O
\emph{eruv} é uma fronteira fluida, uma territorialidade elástica de
apropriação e de pertencimento --- um limite e um limiar em si.

\begin{center}\adforn{49}\end{center}

\noindent{}O \emph{eruv} é um fio de diâmetro milimétrico, transparente. Uma
demarcação e uma delimitação espacial judaica real e imaginária,
concreta e abstrata, visível e invisível, física e ficcional, literária
e literal, legal e comunitária. O \emph{eruv} é uma geografia de
identidade temporária e apropriada. São como ilhas em meio ao oceano
que, antes mesmo de se mostrarem, desaparecem --- para que então possam
nascer em outras localidades e em outros tempos.

Em momentos de diáspora estamos todos nós em busca de um refúgio. Como
menciona o filósofo Denètém Touam Bona\footnote{\textsc{bona}, \emph{op.\,cit.}, p.\,52.}:
``destinados a nos tornar refugiados em nossas próprias terras,
todos estamos em busca de um solo, um ar, uma água que não tenham sido
corrompidos, estamos todos em busca de um fora''\footnote{\emph{Ibidem}.
  Denètém Touam Bona, filho de pai centro-africano e de mãe francesa,
  procura construir pontes e sentidos para uma identidade fronteiriça em
  um mundo que, ainda hoje, torce a ``linha da cor''. Em seus
  projetos, obras e escritos, o autor busca fazer do aquilombamento (a
  arte da fuga e evasão dos escravizados) um objeto filosófico por si
  só, usado para pensar o mundo contemporâneo, as fronteiras, os
  corpos/\,comunidades desviantes, bem como as constantes necessidades de
  fuga --- de territorialidades físicas ou imaginárias.}. Em uma
perspectiva diaspórica, o \emph{eruv} é uma espécie de fora, de outras
temporalidades e territorialidades de agregação e de apropriação que
surgem para além de um espaço (e de um tempo) físico homogêneo e
hegemônico que nos é constantemente apresentado. O \emph{eruv} é uma
fronteira e, assim como todo limite, ``opera a inscrição espacial
de uma coletividade humana em um dado território''\footnote{\emph{Ibidem}, p.
  42.}, simplesmente por instaurar e delimitar este espaço. Como nos
lembra Glória Anzalduá, as fronteiras são esses espaços-territórios de
encontro com o outro --- sejam eles geográficos ou linguísticos.

No livro-ensaio \emph{TAZ}, o filósofo e escritor Hakim Bey sustenta que
as chamadas ``zonas autônomas temporárias'' são táticas
sociopolíticas para criar espaços temporários que escapam às estruturas
formais de controle. Estas outras frequências, que se firmam enquanto
possibilidades de reinventar e reivindicar territórios, são
``operações de guerrilha'' que têm por intuito a libertação de
uma determinada área de terra, de tempo (e de imaginação) por um
deliberado momento.

O \emph{eruv} é um signo de pertencimento, um território de agregação
simbólico e diaspórico, coletivo e comunitário --- um espaço
essencialmente móvel e temporário. Nesta perspectiva, o \emph{eruv} pode
ser compreendido enquanto uma ``zona autônoma temporária'' na
medida em que se fixa, para logo dissolver-se e se refazer em outros
territórios e em outros tempos.

Os espaços, as territorialidades, os locais de pertencimento e de
apropriação --- sejam eles físicos, comunitários ou imaginários --- se
mostram enquanto um conglomerado de tempos e de situações. Nestas
localidades o tempo em si não é linear, ou cíclico/\,elipsoidal, é antes
um ``rizoma temporal''\footnote{\textsc{rago}, \emph{op.\,cit.}, p.\,13.}, ou seja,
um espaço heterotópico em que há o acréscimo de múltiplos e simultâneos
tempos em uma só localidade\footnote{O tempo não é elipsoidal, muito
  menos linear. Trata-se de pensar o tempo enquanto uma multiplicidade:
``coexistência de múltiplas temporalidades em um mesmo momento,
  objeto, situação, relação''. Trata-se, aqui, de pensar o tempo como
  um rizoma temporal. \emph{Ibidem}, p.\,11.}. As heterotopias são alternativas para a
contestação dos espaços existentes, uma possibilidade palpável de
reinventarmos e de darmos novos sentidos às espacialidades físicas,
geográficas, políticas, afetivas e subjetivas, já que ``no seio de
uma heterotopia existe uma ruptura do tempo real''\footnote{\emph{Ibidem}, p.\,14.}.

Porém, ``ao invés de nos deslocarmos para um lugar irreal, para um
espaço e tempo imaginários, seria mais interessante nos voltarmos para o
que nos circunda, estranharmos o que é familiar, percebê-lo
diferentemente e reinventá-lo''\footnote{\textsc{foucault} \emph{apud} ibidem, p.\,14.}.
Os espaços heterotópicos inquietam pois dizem sobre o ``aqui'' e
o ``agora'' --- uma possibilidade de transformar o mundo interior
e exterior, individual e coletivamente. São nestes espaços-tempos
heterotópicos, múltiplos e simultâneos que se dão outros sentidos para
as agregações comunitárias e coletivas.
